%% ==============================================================================
%% content/01-introduction.tex
%% Chapitre 1 : Introduction générale et motivation
%% ==============================================================================
\chapter{Introduction Générale}
\label{chap:introduction}

\section{Contexte et Motivations}
\label{sec:intro-contexte}

La rédaction d'un mémoire universitaire ou d'un ouvrage technique de référence constitue une étape charnière dans l'activité scientifique.
Depuis plusieurs décennies, le système typographique \TeX{} initié par \textcite{knuth1984texbook},
complété par l'environnement logiciel \LaTeX{} conçu par \textcite{lamport1994latex},
s'impose comme la référence incontournable de la publication scientifique rigoureuse.
Cependant, à mesure que l'ampleur des documents augmente et que les exigences de diffusion se diversifient,
la gestion d'un fichier source unique ou d'un préambule monolithique présente des limites criantes.

Dans la pratique courante, les auteurs accumulent souvent des centaines de lignes de configurations hétérogènes,
mélangeant la déclaration des polices,
le réglage des marges,
les redéfinitions de commandes mathématiques
et l'inclusion de paquets tiers incompatibles.
Cette approche engendre une fragilité technique majeure :
la moindre modification d'un style ou le changement de moteur de compilation risque d'altérer l'ensemble de la structure.
L'objectif de cette recherche est de formaliser une architecture logicielle pour \LaTeX{}
fondée sur le principe éprouvé de la séparation des préoccupations (\textit{Separation of Concerns}).

\section{Objectifs et Verrous Scientifiques}
\label{sec:intro-objectifs}

Pour répondre à ces enjeux, nous identifions trois verrous méthodologiques essentiels :

\begin{enumerate}
  \item \textbf{Le découplage sémantique :}
  séparer rigoureusement la logique métier (le texte et les équations)
  de la charte graphique globale et de l'orchestration des paquets.
  \item \textbf{La flexibilité du support de diffusion :}
  permettre une transition transparente entre une version optimisée pour la consultation numérique sur écran
  (liens hypertexte interactifs, marges symétriques)
  et une version destinée à l'impression papier de haute qualité
  (compensation de la reliure \code{BCOR}, liens désaturés, neutralité chromatique).
  \item \textbf{La modularité et la traçabilité collaborative :}
  assurer un découpage granulaire des chapitres et des styles
  facilitant le suivi de version sous Git et minimisant les conflits de fusion.
\end{enumerate}

\begin{callout}[Principe d'architecture]
Un document \LaTeX{} de qualité professionnelle doit être abordé comme un projet logiciel :
le fichier maître ordonnance,
la classe encapsule les règles de composition,
les fichiers d'options adaptent le livrable,
et les fichiers de contenu portent la substance intellectuelle.
\end{callout}

\section{Organisation du Document}
\label{sec:intro-plan}

Le présent document est articulé en quatre chapitres principaux :

\begin{itemize}
  \item Le \cref{chap:etat-art} dresse un panorama critique des approches de mise en page académique,
  en s'appuyant notamment sur les préceptes de \textcite{bringhurst2004elements}
  et sur l'écosystème KOMA-Script documenté par \textcite{kohm2024koma}.
  \item Le \cref{chap:methodologie} formalise l'architecture modulaire retenue,
  le mécanisme des fichiers \code{.clo} et la structuration du flux de travail.
  \item Le \cref{chap:resultats} présente les résultats d'évaluation,
  illustrant la robustesse du modèle face à des contraintes typographiques variées.
  \item Enfin, les annexes et la bibliographie rassemblent les compléments d'implémentation.
\end{itemize}

